\newpage
\section*{PHỤ LỤC}
\phantomsection\addcontentsline{toc}{section}{PHỤ LỤC}

\subsection*{A. Trích đoạn ứng dụng Ryu campus\_switch\_13.py}

\begin{lstlisting}[style=customcode, language=Python, caption={Đồ thị liên kết và tham số thăm dò của khối SDN}]
DIST = (5, 8)
ACCESS = (68, 66, 70, 69)
MGMT_VLAN = 99
CONTROLLER_IP = '10.1.99.10'

# Lien ket: id -> a=(dpid, cong), b=(dpid|'C1'|'C2', cong), w=trong so
LINKS = collections.OrderedDict([
    ('D1-C1', {'a': (5, 'ens9'),  'b': ('C1', 'Et0/2'), 'w': 1}),
    ('D2-C1', {'a': (8, 'ens10'), 'b': ('C1', 'Et1/2'), 'w': 2}),
    ('D1-C2', {'a': (5, 'ens10'), 'b': ('C2', 'Et1/2'), 'w': 1}),
    ('D2-C2', {'a': (8, 'ens9'),  'b': ('C2', 'Et0/2'), 'w': 2}),
    ('D1-D2', {'a': (5, 'ens8'),  'b': (8, 'ens8'),     'w': 2}),
    ('A1-D1', {'a': (68, 'ens4'), 'b': (5, 'ens4'),     'w': 1}),
    # ... A2..A4-D1 (w=1), A1..A4-D2 (w=2)
])
PROBE_ETHERTYPE = 0x88B5
PROBE_INTERVAL = 0.5        # OVS<->OVS
CORE_PROBE_INTERVAL = 1.0   # ARP toi SVI Core
LINK_TIMEOUT = 2.0          # OVS<->OVS: mat probe > 2 s -> chet
CORE_TIMEOUT = 3.5          # OVS<->Core
\end{lstlisting}

\begin{lstlisting}[style=customcode, language=Python, caption={Tính cây dữ liệu tập trung (rút gọn)}]
def compute_data_tree(connected, link_up, frozen=None):
    """Goc = Core-SW1 neu con lien ket song toi no, nguoc lai Core-SW2.
    frozen: Access mat ket noi controller nhung van chay bang flow cu
    -> giu lam la tren lien ket cha cu."""
    usable = {lid: l for lid, l in LINKS.items()
              if link_up.get(lid) and endpoints_ok(lid, connected, frozen)}
    for root in ('C1', 'C2'):
        if not any(link_nodes(lid)[1] == root for lid in usable):
            continue
        # Dijkstra tu goc tren cac canh dung duoc
        parent, done = {}, set()
        heap = [(0, 0, '', str(root), root, None)]
        while heap:
            d, w, lid, _k, n, frm = heapq.heappop(heap)
            if n in done:
                continue
            done.add(n)
            if frm is not None:
                parent[n] = (frm, lid)
            for m, mlid, mw in adj[n]:
                if m not in done:
                    heapq.heappush(heap, (d + mw, mw, mlid, str(m), m, n))
        return root, {p[1] for p in parent.values()}, parent
    return None, set(), {}
\end{lstlisting}


\subsection*{B. Chính sách tập trung trên vSmart}

\begin{lstlisting}[style=customcode, caption={SLA class, AAR, data policy và control policy trên vSmart (rút gọn)}]
policy
 sla-class SLA_REALTIME
  loss 2
  latency 150
  jitter 50
  fallback-best-tunnel
   criteria latency loss
 sla-class SLA_BUSINESS
  loss 5
  latency 300
  jitter 100
  fallback-best-tunnel
   criteria latency loss
 lists
  vpn-list VPN_LAN
   vpn 1
  data-prefix-list SERVER_FARM
   ip-prefix 10.1.90.0/24
  data-prefix-list DMZ
   ip-prefix 10.1.1.0/28
  site-list ALL_EDGES
   site-id 100-500
  site-list BRANCHES
   site-id 200-500
 app-route-policy AAR_LAN
  vpn-list VPN_LAN
   sequence 10
    match
     protocol 1
    action
     count AAR_ICMP
     sla-class SLA_REALTIME fallback-to-best-path
   sequence 40
    match
     destination-data-prefix-list DMZ
     destination-port 443 80
    action
     count AAR_WEB
     sla-class SLA_BUSINESS fallback-to-best-path
 data-policy DP_BRANCH
  vpn-list VPN_LAN
   sequence 10
    match
     destination-data-prefix-list SERVER_FARM
     destination-port 23
     protocol 6
    action drop
     count TELNET_TO_FARM
   sequence 15
    match
     destination-data-prefix-list DMZ
     destination-port 22 23
     protocol 6
    action drop
     count ADMIN_TO_DMZ
   default-action accept
 control-policy PREFER-VE1
  sequence 10
   match route
    originator 10.200.100.1
   action accept
    set
     preference 200
  ! sequence 20/30/40: originator 10.200.200.1,
  ! 10.200.30.1, 10.200.40.1 -> preference 200
  default-action accept
apply-policy
 site-list ALL_EDGES
  control-policy PREFER-VE1 out
  app-route-policy AAR_LAN
 site-list BRANCHES
  data-policy DP_BRANCH from-service
\end{lstlisting}

\subsection*{C. Cấu hình tường lửa chi nhánh}

\begin{lstlisting}[style=customcode, caption={Giao diện, ACL, DHCP, IP SLA và tuyến trên Brand-FW Site 200 (rút gọn)}]
interface GigabitEthernet0/0.60
 vlan 60
 nameif vlan60
 security-level 100
 ip address 10.2.60.1 255.255.255.0
interface GigabitEthernet0/1
 nameif outside1
 security-level 0
 ip address 10.2.1.1 255.255.255.252
access-list SDWAN_IN extended permit ip 10.0.0.0 255.0.0.0 10.2.0.0 255.255.0.0
access-group SDWAN_IN in interface outside1
access-group SDWAN_IN in interface outside2
access-list VLAN99_IN extended permit ip 10.2.99.0 255.255.255.0 10.0.0.0 255.0.0.0
access-list VLAN99_IN extended deny ip any any log warnings
access-group VLAN99_IN in interface vlan99
dhcpd dns 10.1.90.10
dhcpd domain campus.internal
dhcpd address 10.2.60.100-10.2.60.199 vlan60
dhcpd enable vlan60
sla monitor 10
 type echo protocol ipIcmpEcho 10.1.255.1 interface outside1
 timeout 1000
 threshold 1000
 frequency 2
sla monitor schedule 10 life forever start-time now
sla monitor 20
 type echo protocol ipIcmpEcho 9.9.9.9 interface outside2
 timeout 2000
 threshold 2000
 frequency 10
sla monitor schedule 20 life forever start-time now
track 1 rtr 10 reachability
track 2 rtr 20 reachability
! Noi bo: MPLS chinh, Internet du phong
route outside1 10.0.0.0 255.0.0.0 10.2.1.2 1 track 1
route outside2 10.0.0.0 255.0.0.0 10.2.1.6 2
route outside1 10.1.255.1 255.255.255.255 10.2.1.2 1
! Internet: DIA qua vEdge2 chinh, qua MPLS du phong
route outside2 0.0.0.0 0.0.0.0 10.2.1.6 1 track 2
route outside1 0.0.0.0 0.0.0.0 10.2.1.2 2
route outside2 9.9.9.9 255.255.255.255 10.2.1.6 1
\end{lstlisting}

\subsection*{D. Hướng dẫn tái lập hệ thống và phép đo}
\setcounter{table}{0}\renewcommand{\thetable}{D.\arabic{table}}

Phụ lục này tóm tắt các bước dựng lại hệ thống và chạy lại các phép đo trong báo cáo. Toàn bộ topology nằm trong một tệp lab EVE-NG (\texttt{Campus Network SDN SD-WAN.unl}, 74 node, 107 phân đoạn mạng); cấu hình nguồn của từng thiết bị nằm trong thư mục \texttt{configs/} của kho mã, mỗi site một thư mục. Địa chỉ và tham số chính xem ở Bảng~\ref{tab:cac_site}, \ref{tab:vlan_campus_chinh}, \ref{tab:vlan_chi_nhanh} và \ref{tab:system_ip_tloc}.

\textbf{D.1. Thứ tự khởi động.} Các nhóm node được bật theo thứ tự phụ thuộc, nhóm sau chờ nhóm trước sẵn sàng (Bảng~\ref{tab:thu_tu_khoi_dong}). Toàn bộ lab lên đủ sau khoảng 25--30 phút trên máy chủ 48 luồng, 188 GB RAM. Thứ tự này tránh hai lỗi đã gặp: OVS bật trước Controller làm bão broadcast trên VLAN 99, và vEdge bật trước vBond phải chờ thêm một chu kỳ kết nối lại.

\begingroup
\small
\setlength{\tabcolsep}{4pt}
\begin{longtable}{|P{1.0cm}|L{5.6cm}|L{\dimexpr\textwidth-1.0cm-5.6cm-2.2cm-8\tabcolsep-5\arrayrulewidth\relax}|L{2.2cm}|}
    \caption{Thứ tự khởi động các nhóm node}\label{tab:thu_tu_khoi_dong} \\
    \hline
    \textbf{Bước} & \textbf{Node} & \textbf{Lý do} & \textbf{Chờ sau đó} \\
    \hline
    \endfirsthead
    \hline
    \textbf{Bước} & \textbf{Node} & \textbf{Lý do} & \textbf{Chờ sau đó} \\
    \hline
    \endhead
    \hline
    \endfoot
    1 & Router Internet, MPLS; switch Site 900; vBond, vSmart, vManage & Underlay và controller SD-WAN phải có trước mọi Edge & 120 s \\
    \hline
    2 & SwitchServerFarm, Core-SW1/2, SwitchDMZ, SW-WAN-Edge1/2 & Hạ tầng lớp 2/3 và VLAN 99 của Campus chính & 45 s \\
    \hline
    3 & FW-ASAv-Active, rồi FW-ASAv-Standby & Đơn vị Active lên trước để giữ vai trò & 90 s, 45 s \\
    \hline
    4 & vEdge1, vEdge2 Site 100 & Kết nối Campus chính vào fabric & 30 s \\
    \hline
    5 & SDN Controller & Ryu phải nghe cổng 6653 trước khi OVS kết nối & 90 s \\
    \hline
    6 & Dist-SW1/2, rồi Access-SW1--4 & OVS kết nối Controller qua VLAN 99 & 60 s, 30 s \\
    \hline
    7 & DHCP-Server, Web, Mail, Syslog; máy trạm Site 100 & Dịch vụ và máy trạm & 30 s \\
    \hline
    8 & Từng chi nhánh: vEdge1, vEdge2, Brand-FW, switch, máy trạm (Site 200, 300, 400, 500) & Mỗi chi nhánh lên fabric rồi mới cấp DHCP & 30 s mỗi site \\
\end{longtable}
\endgroup

\textbf{D.2. Lệnh kiểm chứng trạng thái.} Sau khi khởi động, trạng thái từng lớp được kiểm bằng các lệnh ở Bảng~\ref{tab:lenh_kiem_chung}; hệ thống chỉ được coi là sẵn sàng khi mọi dòng cho đúng kết quả mong đợi.

\begingroup
\small
\setlength{\tabcolsep}{4pt}
\begin{longtable}{|L{2.6cm}|L{5.4cm}|L{\dimexpr\textwidth-2.6cm-5.4cm-6\tabcolsep-4\arrayrulewidth\relax}|}
    \caption{Lệnh kiểm chứng và kết quả mong đợi}\label{tab:lenh_kiem_chung} \\
    \hline
    \textbf{Thiết bị} & \textbf{Lệnh} & \textbf{Kết quả mong đợi} \\
    \hline
    \endfirsthead
    \hline
    \textbf{Thiết bị} & \textbf{Lệnh} & \textbf{Kết quả mong đợi} \\
    \hline
    \endhead
    \hline
    \endfoot
    SDN Controller & \texttt{curl -s http://127.0.0.1:8080/\allowbreak stats/switches} & Đủ 6 DPID: 5, 8, 66, 68, 69, 70 \\
    \hline
    OVS & \texttt{ovs-vsctl get bridge br0 protocols fail\_mode stp\_enable} & \texttt{OpenFlow13}, \texttt{secure}, \texttt{false} \\
    \hline
    OVS & \texttt{ovs-ofctl -O OpenFlow13 dump-flows br0} & Có flow bảng 0 (kiểm soát) và bảng 1 (chuyển tiếp) \\
    \hline
    Core-SW1/2 & \texttt{show vrrp brief}, \texttt{show ip ospf neighbor} & Core-SW1 Master các VLAN 10--40; láng giềng OSPF \texttt{FULL} \\
    \hline
    FW-ASAv & \texttt{show failover state} & Một đơn vị Active, một Standby Ready \\
    \hline
    vEdge & \texttt{show control connections} & Kết nối \texttt{up} tới vSmart và vManage \\
    \hline
    vEdge & \texttt{show bfd sessions} & Mọi phiên \texttt{up} trên cả hai color \\
    \hline
    vEdge & \texttt{show omp routes vpn 1} & Có prefix LAN của mọi site, trạng thái \texttt{C,I,R} \\
    \hline
    vSmart & \texttt{show policy from-vsmart} (trên vEdge) & Có SLA class, AAR, data policy \\
    \hline
    Brand-FW & \texttt{show track}, \texttt{show route} & Hai đối tượng track \texttt{Up}; tuyến mặc định qua vEdge2 \\
    \hline
    Máy trạm & \texttt{ip dhcp} (VPCS), \texttt{ping} gateway & Nhận địa chỉ trong dải \texttt{.100--.199}; ping thành công \\
\end{longtable}
\endgroup

\textbf{D.3. Công cụ đo.} Các phép đo trong báo cáo được tự động hóa bằng script Python trong kho mã (thư mục \texttt{scripts/} và \texttt{.claude/skills/campus-network-lab/scripts/lab/}). Script đăng nhập máy chủ EVE-NG, mở console thiết bị qua kênh chuyển tiếp của SSH, hoặc chạy đầu đo trong network namespace gắn vào phân đoạn mạng cần đo. Mật khẩu được đọc từ tệp cục bộ lúc chạy, không nằm trong mã. Các script chính: \texttt{sdn\_realcut\_test.py} (cắt liên kết thật trên OVS, đo gián đoạn), \texttt{sdn\_vlan\_test.py} và \texttt{sdn\_policy\_test.py} (thêm VLAN, áp chính sách), \texttt{sdn\_load\_test.py} (tải cao), \texttt{branch\_failover\_test.py} và \texttt{fw\_failover\_test.py} (cắt đường truyền, vEdge, tường lửa), \texttt{icmp\_probe.py} (đầu đo ICMP 5 gói/giây, phân biệt luồng mới và luồng đang chạy), \texttt{site\_join\_test.py}, \texttt{cfgload\_test.py}, \texttt{sdwan\_policy\_time.py}, \texttt{sla\_stats.py}, \texttt{bw\_util.py} và \texttt{mttr\_test.py} (năm tiêu chí vận hành), \texttt{sdwan\_reliability.py} (độ khả dụng từ API vManage).

\textbf{D.4. Cắt liên kết thật.} Mọi phép đo phục hồi cắt liên kết bằng cách hạ giao diện ảo của node trên máy chủ (\texttt{ip link set vunl0\_<node>\_<cổng> down}), giữ một khoảng thời gian cố định rồi bật lại. Thiết bị hai đầu không nhận được sự kiện mất sóng mang, giống đứt cáp ở giữa đường truyền, nên thời gian đo được gồm cả thời gian phát hiện của từng giao thức.
